%%%%%%%%%%  scan idx 622  |  folio 609  %%%%%%%%%%
Il
existe une extension finie \(K\) de \(k(S)\) telle que \((f_K)_*({O}_{X'_K})\) ait toutes ses extensions résiduelles triviales. Si \(S'\) est un schéma intègre fini et plat sur un ouvert \(V\) de \(S\), de corps de fonctions \(k(S')=K\), il existe un ouvert \(V'\) de \(S'\) tel que
\[
  \bigl[(f_{S'})_*({O}_{X'_{S'}})\bigr]_{\red}
  \simeq
  \bigl[(f_{S'})_*({O}_{X_{S'}})\bigr]_{\red}
\]
au-dessus de \(V'\)~; quitte à restreindre \(S\) et \(S'\), on peut supposer \(S=V\) et \(S'=V'\). Comme l'assertion à démontrer est locale pour f.p.q.c., on peut remplacer \(S\) par \(S'\). Enfin on applique 3.1.

\textbf{3.6. Remarque.} Il résulte de même, compte tenu des suites exactes \(\underline{\Pic}_{X_{i-1}/S}\longrightarrow\underline{\Pic}_{X_i/S}\longrightarrow\mathbb V(Q_i)\), que pour tout \(s\in S\)
\[
  \operatorname{coker}\bigl(\underline{\Pic}^0_{X_s/s}\longrightarrow\underline{\Pic}^0_{Y_s/s}\bigr)
\]
est un groupe unipotent.

\section*{4. Démonstration de II* : la partie la plus délicate de la démonstration}

Rappelons l'assertion II*~: le théorème 1.1 lorsque \(X\) et \(Y\) sont intègres sur \(S\) et la suite (*) \({O}_Y\longrightarrow f_*{O}_X\rightrightarrows g_*({O}_{X\times_YX})\) est exacte~; en appliquant le théorème de platitude générique EGA IV 6.9.1 et EGA III 7.8.8, on se ramène à prouver la

\textbf{Proposition 4.1.} \uline{Soient} \(S\) \uline{un schéma noethérien intègre,} \(\lambda:X\to S\) \uline{et} \(\mu:Y\to S\) \uline{des morphismes propres de présentation finie,} \(f:X\to Y\) \uline{un} \(S\)\uline{-morphisme surjectif et} \(g:X\times_YX\to Y\) \uline{le morphisme canonique. On suppose que}
\begin{enumerate}[label=\arabic*)]
\item \(X,Y,X\times_YX\) \uline{sont plats sur} \(S\).
\item \(\lambda_*{O}_X={O}_S\), \(\mu_*{O}_Y={O}_S\) \uline{universellement.}
